\chapter{Architecture}
\label{chap:architecture}

Everything that reads or writes a review goes through the Core API. The dashboard, the MCP server, and the GitHub webhook are clients of that API. The webhook handler does not review code itself: it checks the signature, drops duplicates, enqueues a job, and returns. Indexing and reviewing run in BullMQ workers so a slow model call cannot block the HTTP response.

\begin{figure}[H]
  \centering
  \resizebox{\linewidth}{!}{% System context. Every client reaches data through the Core API.
\begin{tikzpicture}[x=1cm,y=1cm]
  \node[external] (claude) at (1.65,5.15) {Claude Code};
  \node[sys] (mcp) at (1.65,3.45) {MCP server};
  \node[sys] (dash) at (1.65,1.75) {Dashboard\\Next.js};
  \node[external] (gh) at (1.65,0.05) {GitHub};

  \node[sys, text width=2.9cm, minimum width=3.2cm, minimum height=1.25cm] (api) at (6.15,2.6)
    {Core API\\NestJS};

  \node[store] (pg) at (10.35,5.15) {Postgres\\+ pgvector};
  \node[store] (redis) at (13.55,5.15) {Redis\\BullMQ};
  \node[sys] (worker) at (10.35,2.55) {Worker};

  \node[external, text width=2.15cm, minimum width=2.3cm] (ts) at (8.55,0.35) {Tree-sitter};
  \node[external, text width=2.15cm, minimum width=2.3cm] (sg) at (11.15,0.35) {Semgrep};
  \node[external, text width=2.15cm, minimum width=2.3cm] (llm) at (13.75,0.35) {LLM providers};

  \draw[flow] (claude) -- (mcp);
  \draw[flow] (mcp.east) -- node[flabel, above] {REST} (api.north west);
  \draw[flow] (dash.east) -- node[flabel, above] {REST} (api.west);
  \draw[flow] (gh.east) -- node[flabel, below] {webhooks} (api.south west);
  \draw[flow] (api.north east) -- (pg.west);
  \draw[flow] (api.east) -- node[flabel, above, sloped] {enqueue} (redis.west);
  \draw[flow] (redis.south) -- node[flabel, right, sloped] {jobs} (worker.north east);
  \draw[flow] (worker.north) -- (pg.south);
  \draw[flow] (worker.south) -- (sg.north);
  \draw[flow] (worker.south west) -- (ts.north);
  \draw[flow] (worker.south east) -- (llm.north);
\end{tikzpicture}
}
  \caption{Runtime context. The worker posts the finished review back to GitHub. Postgres is the source of truth; Redis carries jobs and live events.}
  \label{fig:architecture}
\end{figure}

\section{Processes}

\begin{description}
  \item[Dashboard.] A Next.js application. Sign-in is Clerk. After sign-in the browser calls the API with the Clerk session.
  \item[Core API.] A NestJS service. It owns REST routes, the webhook endpoint, API-key checks, and the server-sent event streams.
  \item[Worker.] A Node process that consumes the \code{review}, \code{index-repo}, and \code{eval} queues. It is the only process that clones a repository, runs Semgrep, or calls a model during a review.
  \item[MCP server.] A small process that speaks the Model Context Protocol over stdio or Streamable HTTP. It calls the Core API and does not open Postgres or GitHub itself.
  \item[Postgres.] Relational data and the embedding index (pgvector) in one database.
  \item[Redis.] BullMQ queues, rate limiting, and a pub/sub channel for live dashboard updates.
\end{description}

Tree-sitter and Semgrep are tools the worker shells out to or loads in process. Tree-sitter is the WebAssembly build (\code{web-tree-sitter}), so the worker does not need a native compiler. Semgrep is the Semgrep CLI, installed in the worker image. If Semgrep is missing, that stage is skipped and the review continues.

\section{One language, shared types}

The apps and packages are TypeScript, compiled as ES modules. A relative import in a package uses a \code{.js} suffix because the packages emit \code{dist/} with \code{tsc}, and Node resolves that suffix at runtime. Change a package, then build it, before an app will see the change. \code{pnpm dev} builds packages first.

Request and response bodies are Zod schemas in \code{packages/shared}. The API validates with those schemas. The dashboard and the MCP server import the same types. A new field is added once.

Model calls go through \code{packages/llm}. Business code does not import a provider SDK, except that the Claude implementation uses the official Anthropic SDK inside that package. The call is \code{generateJson}: a system prompt, a user prompt, and a Zod schema. The result is checked against the schema. One retry appends the validation error to the prompt. HTTP 429 and 5xx responses back off. If the primary provider fails, the configured fallback provider is tried. Calls are spaced out because the free tiers used in the demo are small.

The provider and model names come from configuration: \code{LLM\_PROVIDER} and \code{LLM\_MODEL}, with the same pair for the fallback. Accepted providers are Gemini, OpenRouter, any OpenAI-compatible host (NVIDIA NIM or Groq, for example), Anthropic, and OpenAI's own API. Embeddings are Gemini or OpenAI. Claude has no embedding model.

\section{Whose keys are used}

Each signed-in user has a row of provider settings. A job runs with the keys of the user it belongs to: the owner of the GitHub installation, or the person who requested an MCP review or an eval run. Keys stored for a user stay in the database and are not sent back to the browser.

Server environment keys are not shared by default. Only the GitHub App admin may use the keys from the environment. \code{SHARED\_LLM\_KEYS=true} shares them with every signed-in user, which is appropriate for a trusted team demo and not for a public deploy.

\section{Queues}

Three queues, named in \code{packages/shared}:

\begin{itemize}
  \item \code{review} carries one review that already has a row in status \code{queued}.
  \item \code{index-repo} carries a full index or an incremental update of one repository.
  \item \code{eval} carries one run of the sample set.
\end{itemize}

Jobs retry with backoff. The webhook's idempotency check (Chapter~\ref{chap:pipeline}) stops a GitHub retry from enqueueing a second review for the same delivery.
